\subsection{子箭图}

设 C 与 C′ 为箭图。若 \(\operatorname{Som}(\mathrm{C}') \subset \operatorname{Som}(\mathrm{C})\) ，\(\operatorname{Fl}(\mathrm{C}') \subset \operatorname{Fl}(\mathrm{C})\) ，且映射 \(o_{C'}\) 与 \(t_{C'}\) 分别与映射 \(o_{C}\) 与 \(t_{C}\) 在 \(\operatorname{Fl}(\mathrm{C}')\) 上相合，则称 C′ 是 C 的子箭图。此外，若 C 中源与靶都属于 \(\operatorname{Som}(\mathrm{C}')\) 的每个箭都是 C′ 的箭，则称 C′ 是 C 的全子箭图。

设 C 为箭图，\((\mathrm{C}_{i})_{i\in\mathrm{I}}\) 为 C 的子箭图族。称族 \((\mathrm{C}_{i})_{i\in\mathrm{I}}\) 的交，并记作 \(\bigcap_{i\in I}C_{i}\) ，这样的 C 的子箭图：其顶点集为 \(\bigcap_{i\in I}\mathrm{Som}(\mathrm{C}_{i})\) ，其箭集为 \(\bigcap_{i\in I}\mathrm{Fl}(\mathrm{C}_{i})\) 。

